
\section{三个不同层次的交叉验证}
外层最终保留用于最后检查，开发内部的交叉验证用于比较模型，内部早停或校准还可能再划分数据。每一层的“训练集”都是相对于该层而言。Pipeline能保证它所包裹的变换在外部交叉验证训练折中拟合，但不能自动理解模型内部再次划分的对象关系。

对分组数据，应把组标识传给需要它的分割器，并检查所有层都遵守任务要求。常见错误包括外层按对象分组、内层却随机拆行；或先用全部开发目标选择特征，再做看似严格的内部五折。发现这样的情况，应重新组织完整流程，而不是只在报告里补一句“采用交叉验证”。

\section{保存指标时保持方向一致}
scikit-learn的交叉验证通常约定分数越大越好，因此有些误差接口以负值返回，如 \nolinkurl{neg_root_mean_squared_error}。打印成RMSE时取一次负号即可；若随后又用最大值选择“最差误差”，容易选反。推荐把结果表列名明确写成 \texttt{rmse}、\texttt{accuracy}等，并在转换处检查一个小例子。

概率列顺序通过 \texttt{classes\_}核对；变换后的列顺序通过相关特征名接口核对。模型序列化时同时保存环境版本、特征清单和预处理，不从不可信来源加载任意序列化文件。不同版本间模型对象的可移植性不能只凭文件扩展名保证。

\section{随书实验的运行顺序}
在工程根目录安装列出的依赖后，运行 \texttt{python scripts/run\_deep\_experiments.py}执行17章深化实验及真实OGLE节选。该脚本按章运行并保存成功、警告和异常；一章失败不会被改名为通过。只重现真实OGLE节选可运行 \texttt{python experiments/project\_ogle.py}。所有这些输入都已包含或可以由明确规则生成。

重生成中文图形需要系统安装中文字体。源码会查找常见字体家族，工程不附带字体文件。已生成的PDF与PNG可直接用于LaTeX编译，即使本地暂未配置Python中文绘图字体也可以阅读和编译教材。
